\chapter{Champs électrique et gravitationnel}\label{ch:g11:electric-gravitational-fields}

Comment la Terre attire-t-elle une pomme qu'elle ne touche jamais~?
Newton lui-même appelait l'action à distance «~une si grande
absurdité~» qu'il refusait de la défendre. La réponse moderne~: une
masse, ou une charge, modifie l'espace autour d'elle --- elle crée un
\emph{\omterm{def:g11:electric-gravitational-fields:field}{champ}} --- et tout ce qui y est placé
répond au \omterm{def:g11:electric-gravitational-fields:field}{champ} en son propre
emplacement. Nous construisons les deux grands exemples côte à côte,
jusqu'à ce que l'analogie devienne une identité de forme.

\section{De l'action à distance au champ}

\begin{definition}[Champ]\label{def:g11:electric-gravitational-fields:field}
Un \emph{champ}\index{champ} assigne à chaque point de l'espace un
vecteur. Une source (une charge, une masse) crée son champ partout
autour d'elle, qu'il y ait ou non quelque chose pour le sentir~; un
objet test ressent alors une \omterm{def:g10:inertia:force}{force}
déterminée par le champ \emph{en son propre emplacement} seul. Cette
image en deux étapes remplace l'action à distance.
\end{definition}

\begin{remark}
Le \omterm{def:g11:electric-gravitational-fields:field}{champ} n'est pas un tour de
comptabilité~: il transporte de l'\omterm{def:g10:energy-conservation:energy}{énergie},
et les changements de la source s'y propagent à vitesse finie --- la
lumière elle-même est un tel
\omterm{def:g11:electric-gravitational-fields:field}{champ} en voyage
(\cref{ch:g12:light-as-wave}). Ici nous étudions des
\omterm{def:g11:electric-gravitational-fields:field}{champs} qui ne changent pas dans le
temps.
\end{remark}

\section{Le champ électrique}

\begin{definition}[Champ électrique]\label{def:g11:electric-gravitational-fields:efield}
Si une petite charge test $q$ placée en un point $M$ ressent une
\omterm{def:g10:inertia:force}{force} électrique $\vect F$, le \emph{champ
électrique}\index{champ électrique} en $M$ est
\[
\vect E = \frac{\vect F}{q}
\qquad \text{(en \unit{N/C}, aussi écrit \unit{V/m})}.
\]
$\vect E$ ne dépend que des sources et de $M$, non de la charge test~:
doubler $q$ double $\vect F$, laissant le quotient inchangé.
\end{definition}

\begin{proposition}[Force sur une charge]\label{prop:g11:electric-gravitational-fields:force}
Une charge $q$ placée là où le
\omterm{def:g11:electric-gravitational-fields:field}{champ} est $\vect E$ ressent la
\omterm{def:g10:inertia:force}{force} $\vect F = q \vect E$~: dans le sens
de $\vect E$ si $q > 0$, opposé à $\vect E$ si $q < 0$, de grandeur
$F = \abs{q}\,E$.
\end{proposition}

\begin{proof}
Réarranger la définition~; la règle de signe est l'algèbre d'un multiple
négatif.
\end{proof}

\begin{example}[Lire un champ]\label{ex:g11:electric-gravitational-fields:reading}
Une charge test $q = \qty{+2.0e-8}{C}$ en $M$ ressent une
\omterm{def:g10:inertia:force}{force} de \qty{6.0e-4}{N} pointant vers l'est.
Le \omterm{def:g11:electric-gravitational-fields:field}{champ} en $M$ est
$E = \num{6.0e-4} / \num{2.0e-8} = \qty{3.0e4}{N/C}$, pointant vers
l'est~; $q' = \qty{-1.0e-8}{C}$ placée au même point ressent
$F = \num{1.0e-8} \times \num{3.0e4} = \qty{3.0e-4}{N}$ --- pointant
vers l'ouest.
\end{example}

\begin{proposition}[Champ d'une charge ponctuelle]\label{prop:g11:electric-gravitational-fields:point}
Une charge ponctuelle $Q$ crée, à la distance $d$, un
\omterm{def:g11:electric-gravitational-fields:field}{champ} le long de la droite joignant
la charge au point --- s'éloignant de $Q$ si $Q > 0$, vers elle si
$Q < 0$ --- de grandeur
\[
E = k\,\frac{\abs{Q}}{d^2},
\qquad k = \qty{8.99e9}{N.m^2/C^2}.
\]
\end{proposition}

\begin{proof}
La loi de Coulomb (\cref{ch:g11:fundamental-interactions}) donne la
\omterm{def:g10:inertia:force}{force} sur une charge test $q$ à la distance
$d$~: grandeur $k \abs{qQ} / d^2$, le long de la droite de jointure,
répulsive pour des signes identiques. Diviser par $q$ et appliquer le
\cref{prop:g11:electric-gravitational-fields:force}.
\end{proof}

\begin{example}[Une sphère chargée]\label{ex:g11:electric-gravitational-fields:sphere}
Une sphère de Van de Graaff porte $Q = \qty{+5.0e-7}{C}$ (une petite
sphère chargée agit comme une charge ponctuelle en son centre). À
$d = \qty{0.50}{m}$~:
\[
E = \num{8.99e9} \times \num{5.0e-7} / 0.50^2 \approx \qty{1.8e4}{N/C},
\]
pointant radialement hors de la sphère. À \qty{1.0}{m} elle est tombée
au quart~: le $1/d^2$ de la \omterm{def:g10:inertia:force}{force} survit
dans le \omterm{def:g11:electric-gravitational-fields:field}{champ}.
\end{example}

\begin{definition}[Lignes de champ]\label{def:g11:electric-gravitational-fields:lines}
Une \emph{ligne de champ}\index{ligne de champ} est une courbe partout
tangente au \omterm{def:g11:electric-gravitational-fields:field}{champ}, une flèche donnant
la direction du \omterm{def:g11:electric-gravitational-fields:field}{champ}. Trois règles~:
\begin{itemize}
  \item elles partent des charges positives et arrivent aux charges
        négatives~;
  \item là où les lignes se serrent le
        \omterm{def:g11:electric-gravitational-fields:field}{champ} est fort~; là où elles
        s'écartent, faible~;
  \item deux lignes ne se croisent jamais~: le
        \omterm{def:g11:electric-gravitational-fields:field}{champ} a une seule direction en
        chaque point.
\end{itemize}
\end{definition}

\begin{omfigure}
\begin{tikzpicture}[scale=0.82]
  \fill[omDef!20] (0,0) circle (0.32);
  \draw[omDef, thick] (0,0) circle (0.32);
  \node[font=\small] at (0,0) {$+$};
  \foreach \a in {0,45,...,315}
    \draw[-Stealth, omThm, thick] (\a:0.32) -- (\a:1.5);
\end{tikzpicture}\qquad\qquad
\begin{tikzpicture}[scale=0.82]
  \fill[omProp!20] (0,0) circle (0.32);
  \draw[omProp, thick] (0,0) circle (0.32);
  \node[font=\small] at (0,0) {$-$};
  \foreach \a in {0,45,...,315}
    \draw[-Stealth, omThm, thick] (\a:1.5) -- (\a:0.32);
\end{tikzpicture}

{\small \omterm{def:g11:electric-gravitational-fields:lines}{Lignes de champ} d'une charge
ponctuelle~: radialement hors de $+$, vers $-$~; le resserrement près
de la charge figure la croissance en $1/d^2$ du
\omterm{def:g11:electric-gravitational-fields:field}{champ}.}
\end{omfigure}

\section{Le champ uniforme d'un condensateur}

\begin{definition}[Condensateur plan]\label{def:g11:electric-gravitational-fields:capacitor}
Deux plaques métalliques parallèles face à face, proches et portant des
charges opposées, forment un \emph{condensateur}\index{condensateur}.
Entre les plaques le \omterm{def:g11:electric-gravitational-fields:field}{champ} est
\emph{uniforme}\index{champ uniforme}~: même grandeur et même direction
partout --- perpendiculaire aux plaques, de la plaque positive vers la
négative.
\end{definition}

\begin{proposition}[Champ entre les plaques]\label{prop:g11:electric-gravitational-fields:uniform}
Si la \omterm{def:g11:circuits-and-power:voltage}{tension} entre les plaques
(\cref{ch:g11:circuits-and-power}) est $U$ et leur écartement est $d$, le
\omterm{def:g11:electric-gravitational-fields:field}{champ} entre elles a pour grandeur
\[
E = \frac{U}{d} \qquad \text{(en \unit{V/m})}.
\]
\end{proposition}
\admitted

\begin{remark}
L'uniformité et $E = U/d$ sont dérivés honnêtement dans le volume de
l'année~1~; la formule explique l'\omterm{def:g10:orders-of-magnitude:unit}{unité}
\unit{V/m}. Des écarts \omterm{def:g11:lenses-and-eye:lens}{minces} font des
\omterm{def:g11:electric-gravitational-fields:field}{champs} forts~: \qty{100}{V} sur un
millimètre font déjà \qty{e5}{V/m}.
\end{remark}

\begin{example}[Diriger un faisceau d'électrons]\label{ex:g11:electric-gravitational-fields:beam}
Dans un \omterm{def:g10:signals-and-waves:oscilloscope}{oscilloscope}, le faisceau passe
entre des plaques avec $U = \qty{100}{V}$ et $d = \qty{5.0}{mm}$~:
$E = 100 / \num{5.0e-3} = \qty{2.0e4}{V/m}$, et chaque électron ressent
$F = \num{1.60e-19} \times \num{2.0e4} = \qty{3.2e-15}{N}$ vers la
plaque positive --- $\num{4e14}$ fois son
\omterm{def:g10:universal-gravitation:weight}{poids} de \qty{8.9e-30}{N}, donc le
faisceau dévie visiblement (la trajectoire courbe est l'affaire du
\cref{ch:g12:newtons-laws}).
\end{example}

\begin{omfigure}
\begin{tikzpicture}[scale=0.9]
  \draw[omDef, very thick] (0,2.2) -- (5.2,2.2)
    node[right, font=\small] {$+$};
  \draw[omProp, very thick] (0,0) -- (5.2,0)
    node[right, font=\small] {$-$};
  \draw[-Stealth, omThm] (0.8,2.05) -- (0.8,0.15);
  \draw[-Stealth, omThm] (1.8,2.05) -- (1.8,0.15);
  \draw[-Stealth, omThm] (2.8,2.05) -- (2.8,0.15)
    node[pos=0.25, right, font=\small] {$\vect E$};
  \draw[-Stealth, omThm] (4.6,2.05) -- (4.6,0.15);
  \fill[omDef] (3.7,1.4) circle (2.5pt)
    node[right=2pt, font=\small] {$q>0$};
  \draw[-Stealth, omMeth, very thick] (3.7,1.4) -- (3.7,0.6)
    node[left, font=\small] {$q\vect E$};
  \draw[Stealth-Stealth, thin] (-0.35,0) -- (-0.35,2.2)
    node[midway, left, font=\small] {$d$};
\end{tikzpicture}

{\small Le \omterm{def:g11:electric-gravitational-fields:capacitor}{champ uniforme} d'un
\omterm{def:g11:electric-gravitational-fields:capacitor}{condensateur}~: lignes parallèles,
également espacées, de la plaque $+$ à la plaque $-$, $E = U/d$~; une
charge positive est poussée le long du
\omterm{def:g11:electric-gravitational-fields:field}{champ}, une négative contre lui.}
\end{omfigure}

\section{Le champ gravitationnel}

\begin{definition}[Champ gravitationnel]\label{def:g11:electric-gravitational-fields:gfield}
Si une masse test $m$ placée en un point ressent une
\omterm{def:g10:universal-gravitation:force}{force gravitationnelle} (un
\omterm{def:g10:universal-gravitation:weight}{poids}) $\vect P$, le \emph{champ
gravitationnel}\index{champ gravitationnel} y est
\[
\vect g = \frac{\vect P}{m} \qquad \text{(en \unit{N/kg})},
\]
et inversement $\vect P = m \vect g$. Un chapitre antérieur a introduit
le nombre $g$~; $\vect g$ ajoute la direction dans laquelle une pierre
lâchée commence à se déplacer.
\end{definition}

\begin{proposition}[Champ de la Terre]\label{prop:g11:electric-gravitational-fields:earth}
À la distance $d$ du centre de la Terre (masse $M$), le
\omterm{def:g11:electric-gravitational-fields:gfield}{champ gravitationnel} pointe vers le
centre et a pour grandeur
\[
g = G\,\frac{M}{d^2},
\qquad G = \qty{6.67e-11}{N.m^2/kg^2}.
\]
\end{proposition}

\begin{proof}
La \omterm{def:g11:fundamental-interactions:four}{gravitation} universelle
(\cref{ch:g10:universal-gravitation}) donne la
\omterm{def:g10:inertia:force}{force} sur une masse test $m$~: grandeur
$G M m / d^2$, vers le centre~; diviser par $m$. À la surface,
$d = R$ et c'est le $g = GM/R^2 \approx \qty{9.81}{N/kg}$ du chapitre
antérieur.
\end{proof}

\begin{example}[Le champ là où vit la Lune]\label{ex:g11:electric-gravitational-fields:moon}
À la distance de la Lune $d = \qty{3.84e8}{m}$~:
\[
g = \num{6.67e-11} \times \num{5.97e24} / (\num{3.84e8})^2
\approx \qty{2.7e-3}{N/kg}.
\]
Le \omterm{def:g11:electric-gravitational-fields:field}{champ} terrestre ne s'arrête
jamais --- il ne fait que s'\omterm{def:g11:lenses-and-eye:lens}{amincir} en $1/d^2$~; c'est le
\omterm{def:g11:electric-gravitational-fields:field}{champ} qui tient la Lune sur son
\omterm{def:g10:universal-gravitation:satellite}{orbite}.
\end{example}

\begin{remark}[Uniforme près de la surface]\label{rem:g11:electric-gravitational-fields:uniformg}
Sur un laboratoire ou même une montagne, $d$ change à peine par rapport
à $R = \qty{6.37e6}{m}$, donc $\vect g$ est
\omterm{def:g11:electric-gravitational-fields:capacitor}{uniforme} à haute précision~: lignes
verticales, \qty{9.81}{N/kg} constant --- le
\omterm{def:g11:electric-gravitational-fields:field}{champ} terrestre ressemble à celui du
\omterm{def:g11:electric-gravitational-fields:capacitor}{condensateur}, tout comme la Terre
ronde paraît plate depuis un jardin.
\end{remark}

\begin{omfigure}
\begin{tikzpicture}[scale=0.9]
  \fill[omExo!15] (0,0) circle (0.85);
  \draw[omExo, thick] (0,0) circle (0.85);
  \node[font=\small] at (0,0) {Terre};
  \foreach \a in {0,60,...,300}
    \draw[-Stealth, omThm, thick] (\a:2.1) -- (\a:0.95);
  \node[omThm, font=\small] at (30:1.9) {$\vect g$};
\end{tikzpicture}

{\small Le \omterm{def:g11:electric-gravitational-fields:gfield}{champ gravitationnel} de la
Terre~: radial, pointant vers le centre, de grandeur $GM/d^2$ --- le
portrait exact du \omterm{def:g11:electric-gravitational-fields:field}{champ} d'une charge
ponctuelle négative.}
\end{omfigure}

\section{Une idée, deux forces}

Les deux \omterm{def:g11:electric-gravitational-fields:field}{champs} de ce chapitre sont la
même mathématique sous deux costumes~:

\begin{center}
\begin{tabular}{lcc}
 & \textbf{électrique} & \textbf{gravitationnel} \\
\hline
source & charge $Q$ & masse $M$ \\
\omterm{def:g11:electric-gravitational-fields:field}{champ} & $\vect E = \vect F / q$ & $\vect g = \vect P / m$ \\
\omterm{def:g10:orders-of-magnitude:unit}{unité} & \unit{N/C} = \unit{V/m} & \unit{N/kg} \\
source ponctuelle & $E = k \abs{Q} / d^2$ & $g = G M / d^2$ \\
constante & $k = \qty{8.99e9}{N.m^2/C^2}$ & $G = \qty{6.67e-11}{N.m^2/kg^2}$ \\
\omterm{def:g10:inertia:force}{force} sur un test & $\vect F = q \vect E$ & $\vect P = m \vect g$ \\
\omterm{def:g11:electric-gravitational-fields:lines}{lignes de champ} & hors de $+$, vers $-$ & toujours vers la masse \\
\end{tabular}
\end{center}

Chaque formule de gauche devient celle de droite sous le dictionnaire
$q \leftrightarrow m$, $\vect E \leftrightarrow \vect g$,
$k \leftrightarrow G$.

\begin{remark}[Où l'analogie se brise]\label{rem:g11:electric-gravitational-fields:noanalogue}
La charge a deux signes~; la masse un seul. Donc les
\omterm{def:g11:electric-gravitational-fields:efield}{champs électriques} peuvent s'annuler~:
les charges redistribuées d'une boîte métallique effacent tout
\omterm{def:g11:electric-gravitational-fields:field}{champ} externe à l'intérieur (une
\emph{cage de Faraday} --- pourquoi un ascenseur tue le
\omterm{def:g10:signals-and-waves:signal}{signal} de votre téléphone), et la matière en
volume, avec $+$ et $-$ équilibrés, est électriquement silencieuse. La
gravité ne peut être ni écrantée ni neutralisée~: chaque
\omterm{def:g10:orders-of-magnitude:unit}{kilogramme} ajoute son attraction --- ce qui
explique pourquoi la \omterm{def:g10:inertia:force}{force} plus faible d'un
facteur $\num{e39}$ (\cref{exo:g11:electric-gravitational-fields:10})
gouverne l'univers aux grandes échelles.
\end{remark}

\section{Exercices}

\begin{exercise}[$\star$]\label{exo:g11:electric-gravitational-fields:1}
Une charge test $q = \qty{+2.0e-6}{C}$ en un point $M$ ressent une
\omterm{def:g10:inertia:force}{force} de \qty{5.0e-3}{N} pointant vers le nord.
Donner le \omterm{def:g11:electric-gravitational-fields:field}{champ} en $M$ (grandeur et
direction), puis la \omterm{def:g10:inertia:force}{force} sur
$q' = \qty{-4.0e-6}{C}$ placée en $M$.
\end{exercise}

\begin{exercise}[$\star$]\label{exo:g11:electric-gravitational-fields:2}
Une charge ponctuelle $Q = \qty{+5.0}{nC}$ est en $O$.
\begin{enumerate}
  \item Calculer le \omterm{def:g11:electric-gravitational-fields:field}{champ} à
        $d = \qty{30}{cm}$~: grandeur, direction.
  \item À quelle distance le
        \omterm{def:g11:electric-gravitational-fields:field}{champ} est-il deux fois plus
        faible~?
\end{enumerate}
\end{exercise}

\begin{exercise}[$\star$]\label{exo:g11:electric-gravitational-fields:3}
Un \omterm{def:g11:electric-gravitational-fields:capacitor}{condensateur} a $U = \qty{600}{V}$
aux bornes de plaques distantes de \qty{3.0}{mm}. Calculer le
\omterm{def:g11:electric-gravitational-fields:field}{champ}, puis la
\omterm{def:g10:inertia:force}{force} sur un électron y
($e = \qty{1.60e-19}{C}$).
\end{exercise}

\begin{exercise}[$\star$]\label{exo:g11:electric-gravitational-fields:4}
Sur Mars, une sonde de \qty{250}{kg} pèse \qty{930}{N}. Calculer le
\omterm{def:g11:electric-gravitational-fields:field}{champ} de surface martien
$g_{\text{Mars}}$, puis le \omterm{def:g10:universal-gravitation:weight}{poids} d'un
astronaute de \qty{80}{kg} là-bas.
\end{exercise}

\begin{exercise}[$\star$]\label{exo:g11:electric-gravitational-fields:5}
Vrai ou faux, avec une raison chacune~: (a)~deux
\omterm{def:g11:electric-gravitational-fields:lines}{lignes de champ} peuvent se croiser là
où le \omterm{def:g11:electric-gravitational-fields:field}{champ} est fort~; (b)~les
\omterm{def:g11:electric-gravitational-fields:lines}{lignes de champ} partent des charges
négatives et arrivent aux positives~; (c)~des
\omterm{def:g11:electric-gravitational-fields:lines}{lignes de champ} resserrées
\omterm{def:g10:signals-and-waves:signal}{signalent} un
\omterm{def:g11:electric-gravitational-fields:field}{champ} fort.
\end{exercise}

\begin{exercise}[$\star\star$]\label{exo:g11:electric-gravitational-fields:6}
La Station spatiale internationale
\omterm{def:g10:universal-gravitation:satellite}{orbite} à l'altitude \qty{420}{km}
($R = \qty{6.37e6}{m}$, $M = \qty{5.97e24}{kg}$).
\begin{enumerate}
  \item Calculer $g$ à la station~; quelle fraction de la valeur de
        surface~?
  \item Les astronautes flottent. Réconcilier cela avec votre réponse en
        une phrase.
\end{enumerate}
\end{exercise}

\begin{exercise}[$\star\star$]\label{exo:g11:electric-gravitational-fields:7}
Quelle grandeur de \omterm{def:g11:electric-gravitational-fields:field}{champ}
équilibrerait le \omterm{def:g10:universal-gravitation:weight}{poids} d'un électron
($m_e = \qty{9.11e-31}{kg}$)~? Comparer avec les \qty{2.0e4}{V/m} de
l'\cref{ex:g11:electric-gravitational-fields:beam}~; conclure sur la
place de la gravité en physique des électrons.
\end{exercise}

\begin{exercise}[$\star\star$]\label{exo:g11:electric-gravitational-fields:8}
Un grain de poussière de masse \qty{2.0e-6}{kg} flotte au repos entre
des plaques horizontales distantes de \qty{2.0}{cm} avec
$U = \qty{5.0}{kV}$ à leurs bornes.
\begin{enumerate}
  \item Calculer le \omterm{def:g11:electric-gravitational-fields:field}{champ}, puis la
        grandeur de la charge du grain.
  \item La charge du grain est positive~: laquelle des plaques est la
        positive~?
  \item Combien de \omterm{def:g11:fundamental-interactions:elementary}{charges
        élémentaires} le grain porte-t-il~?
\end{enumerate}
\end{exercise}

\begin{exercise}[$\star\star$]\label{exo:g11:electric-gravitational-fields:9}
Des charges $+Q$ et $+4Q$ sont distantes de \qty{30}{cm}. Où sur le
segment qui les joint le \omterm{def:g11:electric-gravitational-fields:field}{champ} total
est-il nul~? Expliquer d'abord pourquoi le point doit se trouver entre
les charges, plus près de la plus petite.
\end{exercise}

\begin{exercise}[$\star\star$]\label{exo:g11:electric-gravitational-fields:10}
Dans un \omterm{def:g11:fundamental-interactions:ladder}{atome} d'hydrogène, le proton
($m_p = \qty{1.67e-27}{kg}$, charge $+e$) et l'électron
($m_e = \qty{9.11e-31}{kg}$, charge $-e$) sont distants de
$d = \qty{5.3e-11}{m}$. Calculer les
\omterm{def:g10:universal-gravitation:force}{forces} électrique et gravitationnelle entre
eux, et leur rapport.
\end{exercise}

\begin{exercise}[$\star\star$]\label{exo:g11:electric-gravitational-fields:11}
En un point $M$, une source seule créerait un
\omterm{def:g11:electric-gravitational-fields:field}{champ} de \qty{300}{N/C} pointant vers
l'est~; une seconde, \qty{400}{N/C} pointant vers le nord. Les
\omterm{def:g11:electric-gravitational-fields:field}{champs} s'ajoutent comme des
vecteurs~: donner le \omterm{def:g11:electric-gravitational-fields:field}{champ} total en
$M$ (grandeur et direction), puis la
\omterm{def:g10:inertia:force}{force} sur $q = \qty{-2.0e-6}{C}$ placée
là.
\end{exercise}

\begin{exercise}[$\star\star\star$]\label{exo:g11:electric-gravitational-fields:12}
À quelle altitude le \omterm{def:g11:electric-gravitational-fields:field}{champ} terrestre
est-il tombé à la moitié de sa valeur de surface~? Et à l'altitude
$h = R$, à quelle fraction~?
\end{exercise}

\begin{exercise}[$\star\star\star$]\label{exo:g11:electric-gravitational-fields:13}
Des charges de $\qty{+2.0}{nC}$ sont en $A$ et $B$, distantes de
\qty{6.0}{cm}. $M$ est sur la médiatrice de $[AB]$, à \qty{4.0}{cm} de
son milieu.
\begin{enumerate}
  \item Calculer la distance $AM$, puis la grandeur du
        \omterm{def:g11:electric-gravitational-fields:field}{champ} que chaque charge crée
        en $M$.
  \item Ajouter les deux
        \omterm{def:g11:electric-gravitational-fields:field}{champs} comme vecteurs (utiliser
        la symétrie) et donner le
        \omterm{def:g11:electric-gravitational-fields:field}{champ} total en $M$, grandeur et
        direction.
\end{enumerate}
\end{exercise}

\begin{exercise}[$\star\star\star$]\label{exo:g11:electric-gravitational-fields:14}
Expliquer, en utilisant les deux signes de la charge, comment les
charges libres d'une boîte métallique annulent un
\omterm{def:g11:electric-gravitational-fields:efield}{champ électrique} externe partout à
l'intérieur --- et pourquoi aucun arrangement de masses ne peut faire de
même pour la gravité. Quel ingrédient impossible un bouclier
gravitationnel exigerait-il~?
\end{exercise}

\begin{exercise}[$\star\star\star$]\label{exo:g11:electric-gravitational-fields:15}
Une petite bille ($m = \qty{0.50}{g}$, $q = \qty{+2.0e-7}{C}$) pend à un
fil dans un \omterm{def:g11:electric-gravitational-fields:capacitor}{champ uniforme}
horizontal $E = \qty{1.0e4}{V/m}$. Au repos, le fil fait un angle
$\theta$ avec la verticale.
\begin{enumerate}
  \item Lister les trois \omterm{def:g10:inertia:force}{forces}~; expliquer pourquoi
        $\tan\theta = qE / (mg)$.
  \item Calculer $\theta$, puis la
        \omterm{def:g10:inertia:inventory}{tension} du fil.
\end{enumerate}
\end{exercise}

\section{Problème~: La gouttelette de Millikan}

\begin{problem}[{Problème du week-end --- peser la charge de l'électron~:
une gouttelette d'huile stationnée en l'air entre deux plaques révèle
que la charge vient par pas indivisibles de $\num{1.60e-19}$
\omterm{def:g11:fundamental-interactions:charge}{coulombs}}]
\label{pb:g11:electric-gravitational-fields:1}
En 1909, Robert Millikan pulvérisait des gouttelettes d'huile entre les
plaques horizontales d'un
\omterm{def:g11:electric-gravitational-fields:capacitor}{condensateur}. En réglant la
\omterm{def:g11:circuits-and-power:voltage}{tension} jusqu'à ce qu'une gouttelette
choisie reste immobile --- la \omterm{def:g10:inertia:force}{force} électrique
équilibrant exactement le
\omterm{def:g10:universal-gravitation:weight}{poids} --- il pouvait peser sa charge.
Notre gouttelette a une masse $m = \qty{4.90e-15}{kg}$ (mesurée d'après
sa chute lente champ éteint~; nous la prenons comme donnée)~; les
plaques sont distantes de $d = \qty{8.0}{mm}$, la plaque
\emph{supérieure} est positive~; la charge de la gouttelette est
négative.

\medskip
\noindent\textbf{Partie I --- La scène.}
\begin{enumerate}
  \item Décrire le \omterm{def:g11:electric-gravitational-fields:field}{champ} entre les
        plaques~; que gagne-t-on avec «~
        \omterm{def:g11:electric-gravitational-fields:capacitor}{uniforme}~»~?
  \item Pour $U = \qty{800}{V}$, calculer $E$.
  \item Dans quelle direction la \omterm{def:g10:inertia:force}{force} électrique
        sur la gouttelette~? Pourquoi la plaque \emph{supérieure} devait-elle
        être la positive~?
  \item Calculer le \omterm{def:g10:universal-gravitation:weight}{poids} de la
        gouttelette.
  \item Écrire la condition de repos comme une égalité de deux grandeurs
        de \omterm{def:g10:inertia:force}{force}.
\end{enumerate}

\medskip
\noindent\textbf{Partie II --- L'équilibre.}
\begin{enumerate}[resume]
  \item À partir de là, exprimer $\abs{q}$ en fonction de $m$, $g$, $d$
        et $U$.
  \item La gouttelette reste au repos pour $U = \qty{800}{V}$~: calculer
        $\abs{q}$.
  \item Diviser par $e = \qty{1.60e-19}{C}$. Combien d'électrons en
        excès la gouttelette porte-t-elle~?
  \item La \omterm{def:g11:circuits-and-power:voltage}{tension} est poussée à
        \qty{900}{V}. Dans quelle direction la
        \omterm{def:g10:inertia:force}{force} résultante pointe-t-elle
        maintenant~?
  \item Montrer que la \omterm{def:g11:circuits-and-power:voltage}{tension}
        d'équilibre est $U = mgd/\abs{q}$~: plus la charge est grande,
        plus la \omterm{def:g11:circuits-and-power:voltage}{tension} est petite.
\end{enumerate}

\medskip
\noindent\textbf{Partie III --- Les pas.}
Une rafale de rayonnement peut arracher ou déposer des électrons sur la
gouttelette~; après chaque changement, Millikan retombe la
\omterm{def:g11:circuits-and-power:voltage}{tension} pour l'équilibre.
\begin{enumerate}[resume]
  \item Après une rafale, l'équilibre exige $U = \qty{1200}{V}$. Calculer
        le nouveau $\abs{q}$. La gouttelette a-t-elle gagné ou perdu un
        électron~?
  \item Des équilibres successifs sont trouvés à $U = 1200$, $800$,
        $600$, $480$ et \qty{400}{V}. Calculer les cinq charges.
  \item Montrer que les cinq sont des multiples entiers d'une seule
        quantité~; donner sa valeur.
  \item Entre équilibres consécutifs, de combien la charge change-t-elle~?
        Interpréter.
  \item Expliquer pourquoi cette gouttelette ne peut jamais rester au
        repos à $U = \qty{700}{V}$, si patiemment qu'on règle.
  \item Millikan a répété cela sur des centaines de gouttelettes~: chaque
        charge mesurée était un multiple entier de la même valeur.
        Énoncer la conclusion, et nommer la quantité mesurée.
\end{enumerate}

\medskip
\noindent\textbf{Partie IV --- Pourquoi il n'y a pas de Millikan
gravitationnel.}
\begin{enumerate}[resume]
  \item Dans le dictionnaire de ce chapitre, que jouent les rôles de
        $q$, $E$ et $qE$ dans l'équilibre de la partie~II~?
  \item Une masse au-dessus pourrait-elle soutenir la gouttelette à la
        place~? Calculer le
        \omterm{def:g11:electric-gravitational-fields:field}{champ} d'une sphère de plomb
        de \qty{1000}{kg} à \qty{1.0}{m} au-dessus de la gouttelette, la
        \omterm{def:g10:inertia:force}{force} sur la gouttelette, et comparer
        avec son \omterm{def:g10:universal-gravitation:weight}{poids}.
  \item L'équilibre exige une attraction vers le haut assez forte pour
        combattre toute la Terre. Pourquoi une plaque de charge peut-elle
        la fournir et aucun arrangement de masses ne le peut~?
  \item La \emph{masse} de la gouttelette change aussi par pas (une
        \omterm{def:g11:fundamental-interactions:ladder}{molécule} d'huile~: environ
        \qty{5e-25}{kg}). Comparer les sauts relatifs en $m$ et en
        $\abs{q}$ quand un électron se pose, et conclure~: pourquoi cet
        équilibre voit-il les
        \omterm{def:g11:fundamental-interactions:ladder}{atomes} de l'électricité mais
        pas les \omterm{def:g11:fundamental-interactions:ladder}{molécules} d'huile~?
        Chute~: la charge est quantifiée, et son
        \omterm{def:g11:fundamental-interactions:ladder}{atome} est
        $e = \qty{1.60e-19}{C}$.
\end{enumerate}
\end{problem}
